\chapter{Your First Program: Talking to the Computer}

\section{Introduction: Real-Life Hook}

Students already “talk” to computers without thinking of it as programming:  

\begin{itemize}
    \item Chat apps: sending messages to friends  
    \item Voice assistants: asking Google or Siri for the weather  
    \item Smart devices: telling Alexa to play music  
\end{itemize}

In all these cases, the computer responds to instructions. Programming is just giving instructions in a language the computer understands.

\section{Concept: What Is a Program?}

\begin{definition}{Program}
A \textbf{program} is a set of instructions that tells the computer what to do.  
Every app, game, or website is controlled by programs.
\end{definition}

\subsection{Talking to the Computer}

Python allows us to write programs that “talk” to the computer.  
The simplest interaction is showing messages on the screen.

\begin{definition}{Print Statement}
\texttt{print()} is a Python command that displays messages to the user.
\end{definition}

\subsection{Example: Print a Message}

\begin{verbatim}
print("Hello, world!")
\end{verbatim}

Output on screen:

\begin{verbatim}
Hello, world!
\end{verbatim}

\section{Comments: Talking to Humans, Not Computers}

Sometimes programmers write notes inside programs to explain the code.  
These are called \textbf{comments}.

\begin{verbatim}
# This program prints a greeting
print("Hello, world!")
\end{verbatim}

The computer ignores comments—they are only for humans.

\section{Diagram: How Your First Program Works}

\begin{center}
\begin{tikzpicture}[node distance=1.5cm]
\node (start) [draw, ellipse, fill=green!20] {Start Program};
\node (instructions) [draw, rectangle, fill=yellow!20, below of=start] {Execute Instructions};
\node (print) [draw, rectangle, fill=orange!20, below of=instructions] {Display Message with print()};
\node (end) [draw, ellipse, fill=green!20, below of=print] {End Program};
\draw[->] (start) -- (instructions);
\draw[->] (instructions) -- (print);
\draw[->] (print) -- (end);
\end{tikzpicture}
\end{center}

\section{Project: Introduce Yourself}

Write a Python program that prints your name, age, and favorite hobby.  

\begin{verbatim}
# My first introduction program
print("Hello! My name is Amina.")
print("I am 15 years old.")
print("My favorite hobby is painting.")
\end{verbatim}

Output:

\begin{verbatim}
Hello! My name is Amina.
I am 15 years old.
My favorite hobby is painting.
\end{verbatim}

\section{Key Concepts}

\begin{itemize}
    \item A program is a set of instructions.
    \item Computers follow instructions exactly.
    \item \texttt{print()} displays messages to users.
    \item Comments (\#) are ignored by the computer and used to explain code.
    \item Python is an easy first language for learning programming.
\end{itemize}

\section{End-of-Chapter Exercises}

\begin{exercise}{Very Short Questions}
\begin{enumerate}
    \item What is a program?
    \item What does \texttt{print()} do in Python?
    \item Why do we use comments in a program?
\end{enumerate}
\end{exercise}

\begin{exercise}{Short Answer Questions}
\begin{enumerate}
    \item Write a Python statement to print your favorite food.
    \item Why is Python suitable for beginners?
    \item Explain the difference between code and comments.
\end{enumerate}
\end{exercise}

\begin{exercise}{Activity-Based Questions}
\begin{enumerate}
    \item Create a program that introduces your best friend.
    \item Modify your introduction program to include your favorite subject and sport.
    \item Ask a classmate to run your program and observe the output.
\end{enumerate}
\end{exercise}

\section{Teacher Notes}

\marginnote{
\textbf{Learning Objectives:} 
\begin{itemize}
    \item Students understand what a program is.
    \item Students learn to display messages using Python.
    \item Students understand the role of comments.
\end{itemize}

\textbf{Teaching Tips:} 
\begin{itemize}
    \item Use chat apps and voice assistants as hooks.
    \item Demonstrate print() statements live.
    \item Encourage students to modify messages for fun.
\end{itemize}

\textbf{Assessment:} 
\begin{itemize}
    \item Check if students’ programs run correctly.
    \item Observe understanding of comments and messages.
\end{itemize}
}